% LSPV-5 REPORT::ARC-SOURCE-RETENTION-QUANTIFIER-FIREWALL229::2026-10-30T17:48:00-07:00
\documentclass[11pt]{article}
\usepackage[margin=.9in]{geometry}
\usepackage{amsmath,amssymb,amsthm,booktabs,array,hyperref,xurl}
\newtheorem{theorem}{Theorem}
\newtheorem{remark}{Remark}
\newcommand{\EXACT}{\textbf{PROVED/EXACT}}
\newcommand{\KNOWN}{\textbf{KNOWN THEOREM}}
\newcommand{\COMP}{\textbf{COMPUTATIONAL}}
\newcommand{\OPEN}{\textbf{OPEN}}
\newcommand{\RET}{\textbf{FALSIFIED/RETIRED}}
\title{LSPV-5 Independent Referee: Source-Retention Quantifier Firewall\\
for Reusable Summaries and ARC/FLMV Contraction}
\author{LSPV-5 independent red team and synthesis}
\date{30 October 2026, 17:48 America/Los\_Angeles}
\begin{document}
\maketitle

\begin{abstract}
The preceding run constructed, for every boundary function
$f:\{0,1\}^b\to\{0,1\}$, a canonical proper 3-CNF $A_f(X,Y)$ of
$N_b=\Theta(b^2 2^b)$ bits and proper point continuations $B_a$ such that
$\operatorname{SAT}(A_f\wedge B_a)=f(a)$.  It proved that a detached reusable
summary prepared before $a$ is known and exact for all $a$ needs at least
$2^b=N_b^{1-o(1)}$ bits.  This run proves the matching interface escape:
when the original canonical source remains indexed read-only and $B_a$ is
known, a fixed controller reads the first literal word of block $a$ and
returns $f(a)$.  It uses one source-word probe, $O(b)$ address/work bits and,
on the native bit-addressed model, $O(b)$ addressed source-bit reads plus
polynomial-in-$b$ address computation.  The two statements coexist because
their prefixes are respectively one state for all future queries and one
targeted probe after the query is known.  Thus a detached-summary lower bound
cannot be promoted to a SAT runtime or space lower bound while source retention
is legal.  Full syntax recognition on an unpromised source costs a charged
scan, and noncanonical inputs require an exact fallback.  No all-input
hereditary $A\mapsto\Phi(A)$ contractor is constructed.  The actual FLMV
forbidden region and P versus NP remain \OPEN.
\end{abstract}

\section{Claim card, missing lemma, and falsifiers}
\paragraph{Single missing lemma tested.}
Does the near-linear detached reusable-summary floor from
ARC-REUSABLE-BOUNDARY-SUMMARY-FIREWALL228 imply that a source-retaining
adaptive residual-SAT controller must spend the same amount of writable space
or time on the realized continuation?

\paragraph{Exact prefix and model.}
The lower-bound interface is
\[
 \exists P,Q\ \forall f:\quad S_f=P(A_f),\quad
 \forall a:\ Q(S_f,B_a)=f(a),
\]
where $Q$ cannot reread $A_f$.  The falsifying retained-source interface is
\[
 \exists R\ \forall f\ \forall a:\quad
 R^{A_f}(B_a)=f(a),
\]
where $R$ receives $a$ before choosing an input address.  The superscript is
indexed read-only access to the same explicit $A_f$, not semantic advice.
The source-word accounting below is first stated in a fixed-width word layout;
native indexed-bit costs are then charged separately.

Four mechanisms were attacked.
\begin{enumerate}
\item A detached summary floor forces the same writable-space floor with the
source retained.  \RET.
\item The controller must prepare answers for all $2^b$ counterfactual
continuations although one execution receives only one $B_a$.  \RET.
\item One word probe may be treated as one native FLMV step.  Rejected as a
model mismatch; the literal word contains $\Theta(b)$ bits.
\item The easy promise-family accessor is already an all-input ARC successor.
Rejected: recognizing the promise and handling nonmembers remain charged, and
a parent-size fallback defeats the desired exponent recurrence.
\end{enumerate}

\section{Canonical family recalled}
For each $a\in\{0,1\}^b$, let
\[
 D_a=\bigvee_{i:a_i=0}x_i\ \vee
     \bigvee_{i:a_i=1}\neg x_i,
\]
which is false exactly at $a$.  The canonical family allocates one fixed block
of $b-2$ proper ternary clauses for every $a$.  If $f(a)=0$, the block is the
standard existential implication-chain encoding of $D_a$; if $f(a)=1$, the
block consists of $b-2$ independent positive dummy clauses on its reserved
private variables.  Blocks have disjoint private ranges and fixed order.
Consequently
\[
 \exists Y\,A_f(x,Y)=\bigwedge_{a:f(a)=0}D_a(x)=f(x).
\]
There are $(b-2)2^b$ clauses.  With fixed-width binary variable identifiers,
the explicit source length is $N_b=\Theta(b^2 2^b)$.

The proper 3-CNF $B_a$ pins every $x_i$ using four clauses and two fresh
variables per pin.  Thus $A_f\wedge B_a$ is satisfiable exactly when $f(a)=1$.
Parsing canonical $B_a$ costs its physical length, $O(b^2)$ bits in this
encoding, and stores the $b$-bit address $a$.  No satisfying assignment or
answer bit is supplied.

\section{Exact retained-source accessor}
\begin{theorem}[Canonical source-retention quantifier separation]
For every $b\ge3$ and every Boolean $f$ in the canonical family above:
\begin{enumerate}
\item every deterministic detached fixed-length summary exact for all point
continuations has worst-case length at least $2^b$ bits;
\item after canonical $B_a$ is available, a fixed deterministic controller
with indexed read-only access to $A_f$ decides
$\operatorname{SAT}(A_f\wedge B_a)$ by reading exactly one literal word of
$A_f$, using $O(b)$ work/address bits;
\item on the native indexed bit-input realization, the same test reads the
$\Theta(b)$ bits of that literal word and uses polynomial-in-$b$ address time,
not one unit-cost native step.
\end{enumerate}
Hence item 1 alone supplies neither a time nor a writable-space lower bound for
source-retaining or joint-input SAT algorithms.
\end{theorem}

\begin{proof}
Item 1 is the deterministic one-way INDEX argument already proved in report228:
the response vector over all $a$ is exactly the truth table of $f$, so distinct
$f$ must have distinct detached summaries.

For item 2, the first clause of block $a$ is at clause address
\[
 p(a)=a(b-2).
\]
When $f(a)=0$, the block is the exclusion chain and its first literal is
$x_1$ or $\neg x_1$, whose identifier has absolute value at most $b$.  When
$f(a)=1$, the block is dummy and its first literal belongs to the reserved
private range, so its identifier is greater than $b$.  Reading that literal
word and testing the identifier therefore returns $f(a)$ exactly.  The
address $p(a)$ and threshold $b$ require $O(b)$ bits because
$p(a)<b2^b$.

For item 3, a variable identifier in this layout has $\Theta(b)$ bits.
An indexed bit machine must address and read those bits, plus compute their
physical offsets; this is polynomial in $b=\Theta(\log N_b)$ and uses $O(b)$
writable bits.  Calling the complete word access one native step would silently
change models.  The conclusion follows because items 1 and 2 concern different
information interfaces and quantifier orders.
\end{proof}

\begin{remark}[Promise recognition and total correctness]
The theorem is exact on the declared canonical promise family.  A total
all-input controller may verify the entire layout with an $O(N_b)$ source scan
and then use the accessor.  If verification fails, it must invoke another exact
procedure.  A full-parent call to arbitrary solver $A$ on that fallback branch
retains the unknown parent exponent $K$ and does not establish
$K'\le\max\{\alpha,\gamma+\lambda K\}$.  Thus the promise recognizer is a
useful easy-family terminal, not the missing hereditary ARC compiler.
\end{remark}

\section{Batch, cache, and sequential-source accounting}
For an execution receiving $q$ point continuations, the direct controller uses
$q$ literal-word probes.  Caching one answer per distinct observed query uses
$r$ probes, $r$ answer bits, and at most $br$ address bits, where $r\le q$ is
the number of distinct realized points.  Nothing requires materializing the
remaining $2^b-r$ counterfactual answers.  At $b=16,q=128$, the reproducible
ledgers give:
\[
\begin{array}{lrrrr}
\toprule
\text{query pattern}&r&\text{no-cache probes}&\text{cached probes}&
\text{cached answer/address bits}\\
\midrule
\text{one point repeated}&1&128&1&1+16\\
\text{128 distinct points}&128&128&128&128+2048\\
\text{eight points repeated}&8&128&8&8+128\\
\bottomrule
\end{array}
\]
The detached all-query floor is $2^{16}=65536$ bits, but it belongs to the
stronger reusable interface.

On a sequential read-only input, the direct implementation may scan to block
$a$; the last block begins after
$3(b-2)(2^b-1)$ literal words.  This is an algorithm-specific upper cost, not
an all-algorithm lower bound.  Conversely, the indexed-source result is not
borrowed into a sequential-tape theorem.  Physical source access, index-tape
updates, output, and peak live state remain model-specific.

\section{Finite exhaustive audit and hostile controls}
The producer exhaustively enumerated all $256$ boundary functions for $b=3$
and all $65536$ functions for $b=4$.  For each function it built the complete
canonical proper 3-CNF, performed a complete syntax scan, and checked every
point continuation by the one-word accessor.  The totals were:
\[
\begin{array}{rrrrr}
\toprule
b&\text{functions}&\text{point queries}&\text{word probes}&
\text{recognizer literal comparisons}\\
\midrule
3&256&2048&2048&6144\\
4&65536&1048576&1048576&6291456\\
\bottomrule
\end{array}
\]
All $1{,}050{,}624$ point answers were correct.  An independently implemented
raw-output checker passed 42 arithmetic, opposite-label, batch, and source
control assertions.

The same canonical layout was tested on a unique-boundary-model SAT source and
the zero-boundary-function UNSAT source.  At point $a=15$, the SAT source's
addressed literal was private variable $95$ and the controller returned $1$;
the UNSAT source's addressed literal was $-x_1$ and it returned $0$.  The same
UP--DPLL backend verified both labels.  Original-clause controls record every
signed occurrence's local all-literals-true $S_j$ target, inverted $U_j$,
local $U_j$ and inverted $S_j$, all SS/SU/US/UU pair compatibilities, coherent
$0^n/1^n$ endpoints and intermediate flip histograms.  They are adversarial
controls, not SAT oracles.

The producer ran in 4.1 seconds and the independent checker in under one
second in this environment; both respect the 30-second chunk rule.  These
wall times and selected probes are not calibrated FLMV steps.  SHA-256:
\begin{itemize}
\item producer: \texttt{2a08cb130b41b30701b985f48ec9973ba1044cfd1f7d1c7794a62d99baa04e98};
\item checker: \texttt{d3a713e5cebcb94a1f7bcd6bfaff4b5f69abd765625cedf572b934cce1b695fa};
\item raw JSON: \texttt{7b378f8df61fcec009aafb969b0c6cd92ff699847737b08f3e07036d41ddd077};
\item checker JSON: \texttt{7eab7cd5a707411333fa84d11347daac5e492deea0ec06f35e565099a5cab82c}.
\end{itemize}
The readback-targeted GitHub revisions are producer commit
\texttt{0db1b35d4f3fbaabb0aea46cc39d0a3cc13d9e0a}, checker commit
\texttt{d970dcf7fd743520e7e1ba1a24fc750c919822e5}, compressed exact-raw commit
\texttt{f84af738b03d5b6bf558bd80cb68d438dd6b8337}, and independent-result commit
\texttt{430789175e05af1b73fa7a8ef95979ce99743272}.  The compressed raw file is
base64 text of gzip-compressed JSON; its encoded-file SHA-256 is
\texttt{8621bb9d70e173db4574a10f4f6301c28d68072e9b81446a45caef3c848ed34e}.

\section{Literature and implementation audit}
Miltersen--Nisan--Safra--Wigderson formalize the connection between static
data structures and asymmetric communication.  It supports the distinction
used here: a one-way all-query message and an online random-access data
structure are different models, with storage/query tradeoffs rather than an
automatic transfer of one lower bound to the other.  The exact $2^b$ detached
summary floor is the elementary deterministic INDEX row-count argument.
Darwiche--Marquis distinguish offline compilation from online queries; their
map does not assert every CNF has a small compiled form.

Actual code was inspected rather than inferred from a summary.  The project
producer at commit \texttt{45da91310af4ecfb935f60c638878cc1aaf41daa}
implements the canonical boundary family used here.  The external DPMC
project was checked at commit
\texttt{a0ee32a652cfd5251b33d257ee2593120c3cd20b}; its project-join dynamic
programming is a supplied/discovered structural architecture, not an
all-input power-short ARC contractor.  No novelty claim is made for the INDEX
or data-structure distinction.

\paragraph{Cross-domain classification.}
Static data structures and asymmetric communication are a
\textbf{THEOREM TRANSFER} only for the explicitly matched preprocessing/query
interface.  Treating their tradeoff as an unrestricted SAT runtime lower bound
is a \textbf{FAILED TRANSFER}.  Knowledge compilation is an
\textbf{ANALOGY/CANDIDATE} for reusable residual summaries until acquisition,
representation, and arbitrary-input coverage are proved and charged.

\section{Corpus reconciliation}
\paragraph{CORPUS REVIEWED.}
The 09:48 governance redirect, UNIVERSAL-PARAMETERIZED-FLMV-CONTRACT1,
AFFINE-SMALLBIAS-DESCENT-GATE2, CONSTRUCTION-STRENGTHEN17,
ARC-LINEAR22-SATURATION1, root \texttt{main.tex}, AGENTS, CODEX, research
skill/protocol, issue \#3, report228, and the newest LSPV-1 through LSPV-4
cycle were inspected from GitHub.  The latest peers prove guarded contextual
clause export, safe transactional aborted elimination, restricted cofactor
mass accounting, and a restricted assignment-state information floor.  None
constructs the arbitrary-input contractor.

\paragraph{DUPLICATION / CONTRADICTION AUDIT.}
Root main already stated in a general detached-summary paragraph that a
bit-vector control uses zero summary bits and one source probe.  The new work
does not claim that observation as novel; it realizes it inside the exact
proper-3CNF family of report228, proves the address decoder, and charges
word/bit-model, recognition, query, and batch costs.  LSPV-4 independently
notes that paid on-demand source reads escape a retained occurrence index; the
present theorem supplies a minimal exact example.  There is no contradiction
with its restricted assignment-state floor because the source cannot recover
an internally generated pin history, while it does contain the static $f(a)$
bit in this canonical family.

\paragraph{WHAT CHANGED RELATIVE TO ROOT MAIN.}
The report228 summary-size theorem remains exact.  Its runtime-transfer
firewall is strengthened from an abstract bit-vector example to the same
proper-3CNF family: near-linear detached summary and polylogarithmic targeted
retained-source access coexist.

\paragraph{ADOPTED / REJECTED / RETIRED PRIOR COMPONENTS.}
Adopted: report228 construction, LSPV-1 source authentication, LSPV-2 safe
rollback, LSPV-3 explicit branch-mass accounting, LSPV-4 same-machine FLMV
discipline.  Rejected: free best-policy selection, all-query counterfactual
storage charges on a one-query execution, unit-cost word access, timeout NO,
and parent-scale fallback as a contraction.  Retired only: promotion of the
detached reusable-summary floor to an unrestricted source-retaining runtime or
space lower bound.

\paragraph{STALE OR OVERSTATED CLAIMS FOUND.}
The newest LSPV-4 report is substantive and has a valid final sentinel, but its
opening line omits the mandated word \texttt{REPORT}; a format-only recovery
request was posted without reconstructing or altering its science.  No
scientific promotion was based on that marker defect.  The affine small-bias
theorem remains restricted affine MCSP semantics and cannot be promoted to SAT
runtime.

\section{Status ledger and next lemma}
\EXACT: the canonical source-retention quantifier theorem and the explicit
address decoder.  \COMP: exhaustive $b=3,4$ audit and matched SAT/UNSAT source
controls.  \RET: using the detached summary floor as a general lower bound on
source-retaining or joint-input SAT computation.  \KNOWN: FLMV's actual
same-machine simultaneous time--space forbidden region.  \OPEN: a uniform
hereditary all-input ordinary-SAT $A\mapsto\Phi(A)$ with all source access,
policy discovery, failed trials, physical child output, guest simulation and
peak retained space charged.  P versus NP remains OPEN.

\paragraph{Next specific lemma.}
For one proposed ARC child interface, formalize a \emph{realized-query access
tariff}: bound, on every execution, the union of source words queried across
all children plus address-generation, repeated-query, ancestor-retention and
fallback costs.  Prove that the tariff composes hereditarily on both SAT and
UNSAT sources without charging all counterfactual queries, and determine
whether it fits the strict $(241/200,11642479/23232400)$ FLMV box.  A pointwise
one-probe family is insufficient.

\begin{thebibliography}{9}
\bibitem{MNSSW} P. B. Miltersen, N. Nisan, S. Safra and A. Wigderson,
\emph{On Data Structures and Asymmetric Communication Complexity}, STOC 1995,
pp. 103--111, DOI 10.1145/225058.225093; JCSS 57 (1998), 37--49.
\bibitem{DM} A. Darwiche and P. Marquis, \emph{A Knowledge Compilation Map},
JAIR 17 (2002), 229--264, DOI 10.1613/JAIR.989.
\bibitem{DPMC} J. Dudek, V. H. N. Phan and M. Y. Vardi,
\emph{DPMC: Weighted Model Counting by Dynamic Programming on Project-Join
Trees}, CP 2020; code \url{https://github.com/vardigroup/DPMC}.
\bibitem{FLMV} L. Fortnow, R. Lipton, D. van Melkebeek and A. Viglas,
\emph{Time-Space Lower Bounds for Satisfiability}, JACM 52 (2005), 835--865,
DOI 10.1145/1101821.1101822.
\end{thebibliography}
\end{document}
% END_OF_VERIFIED_REPORT::LSPV-5::ARC-SOURCE-RETENTION-QUANTIFIER-FIREWALL229
